\documentclass[11pt]{article}
\usepackage[margin=1in]{geometry}
\usepackage{booktabs}
\usepackage{graphicx}
\usepackage{longtable}
\usepackage[T1]{fontenc}
\usepackage[hidelinks]{hyperref}
% Keep typographic protrusion while avoiding pdfTeX font expansion on systems
% where the available T1 fonts are not scalable.
\usepackage[expansion=false]{microtype}

\title{Graduate Student Engagement with a Collaboratively Maintained\\
TA Wiki: A Descriptive Departmental Survey Evaluation}
\author{Antonio Aguirre \and Marcela Alfaro-Cordoba \and Andrew Le}
\date{July 2026}

\begin{document}
\maketitle

\begin{abstract}
This descriptive departmental survey evaluation examines graduate-student reports of
awareness, use, contribution barriers, and perceived value of a collaboratively
maintained teaching-assistant (TA) Wiki. A manifest-bound, disclosure-controlled
analysis was used to summarize structured survey responses. This report presents
only the aggregate result permitted by the pre-specified disclosure policy; all
other planned structured summaries were withheld rather than selectively
reported. The released result concerns respondents' assessment of the Wiki as a
departmental resource. Findings describe this implementation and should not be
interpreted as causal effects or population estimates.
\end{abstract}

\section{Study context and purpose}

This study describes how respondents in one departmental implementation
reported awareness of, use of, and possible engagement with a collaboratively
maintained TA Wiki. The study is descriptive: it does not identify the effect of
the Wiki, GitHub hosting, editathon participation, or any other measured
factor. Its practical aim is to document what can be reported transparently
while protecting respondents in a small voluntary survey.

\section{Methods}

\subsection{Design, setting, and participants}

This was a voluntary departmental survey evaluation conducted under an
institutional exemption determination for the survey-only project. Recruitment
used general announcements to eligible graduate students. The analytic cohort
was defined from the approved consent and eligibility responses in the frozen
live export. The invitation denominator was not verified for reporting, so no
response rate is reported.

\subsection{Instrument and reporting approach}

The analysis covered structured items about background and teaching context,
resource use, Wiki awareness and use, contribution, GitHub, editathon
experience, and perceived value. Optional open-text fields were outside the
automated quantitative pipeline. The reporting structure was informed by
survey-reporting guidance such as CROSS and, where relevant to electronic
questionnaire reporting, CHERRIES \cite{cross,cherries}.

\subsection{Data handling and analysis}

The finalized source export, timestamps, row-level derivatives, and open text
were held only in external restricted storage. Raffle and contact records
remained separate from the analytic source. A manifest-bound workflow validated
the immutable source, wrote a minimal cohort ledger, re-opened the frozen source
only within restricted storage for controlled transformation, and generated
restricted internal outputs. No missing values were imputed.

Structured items were summarized descriptively with item-specific denominators,
counts, and percentages. Category order followed the controlled codebook.
Pre-specified transformation rules handled verified format anomalies, and
incompatible single-choice values were excluded from that item's valid
denominator. The direct contribution response was primary; a conditional
alternative was retained only as an internal sensitivity analysis. Any
cross-tabulations were exploratory, internal-only, and not interpreted as
causal, inferential, or population-generalizable.

\subsection{Disclosure controls and qualitative scope}

The release policy suppressed respondent-level records, timestamps, open text,
and every structured table containing a positive cell below the approved
minimum-cell threshold. Consequently, a withheld table is a disclosure decision
rather than evidence of no pattern. Qualitative material was deferred: no
coded themes, paraphrases, or quotations are reported.

\section{Results}

The disclosure-approved candidate retained one structured summary, presented in
Table~\ref{tab:structured-survey-summaries}. The remaining pre-specified
structured tables were fully withheld under the release policy. Accordingly,
this report makes no item-level claims about awareness, use, contribution
barriers, GitHub, editathon experience, or resource preferences beyond what is
shown in the released table.

\textit{No disclosure-approved numerical result artifact is currently available
for this manuscript. A historic restricted structured run is not a current
release source. This placeholder must be replaced only by generated content
bound to an approved release manifest; numerical results must not be entered
manually.}


\section{Discussion}

The released result concerns perceived departmental value only. Its distribution
indicates a favorable assessment of the TA Wiki among the analytic cohort, while
the disclosure boundary appropriately prevents stronger conclusions about other
survey domains. This is useful evidence for maintaining a departmental resource,
but it does not demonstrate that the Wiki changed teaching practice, increased
engagement, or would yield the same pattern outside this setting.

The report's narrow scope is itself an important limitation. Participation was
voluntary, the invitation denominator was unavailable for a response-rate
calculation, and most planned summaries could not be shared without exposing
small cells. The analysis therefore prioritizes transparent withholding over
apparently comprehensive but unsafe reporting. It also does not include
qualitative findings, which would require a separate manual coding and
disclosure-review process.

\section{Conclusion}

Within this departmental survey evaluation, the disclosure-approved evidence supports
a limited descriptive conclusion about the TA Wiki's perceived value. Future
work should use a larger or differently designed evaluation if the team wishes
to study awareness, adoption, barriers, or outreach effectiveness in greater
detail. Any such work should retain the separation of survey data from raffle
records and apply disclosure controls before sharing results.

\section*{Ethics and data availability}

The survey-only project received an institutional exemption determination. Code,
controlled metadata, and synthetic tests may be shared. Survey data, timestamps,
respondent-level derivatives, open text, and restricted governance records
remain unavailable. The numerical table in this manuscript is generated only
from the approved aggregate release candidate.

\begin{thebibliography}{9}

\bibitem{cross}
Sharma A, Minh Duc NT, Luu Lam Thang T, et al.
A Consensus-Based Checklist for Reporting of Survey Studies (CROSS).
\textit{Journal of General Internal Medicine}. 2021;36(10):3179--3187.
doi:\href{https://doi.org/10.1007/s11606-021-06737-1}{10.1007/s11606-021-06737-1}.

\bibitem{cherries}
Eysenbach G.
Improving the Quality of Web Surveys: the Checklist for Reporting Results of
Internet E-Surveys (CHERRIES).
\textit{Journal of Medical Internet Research}. 2004;6(3):e34.
doi:\href{https://doi.org/10.2196/jmir.6.3.e34}{10.2196/jmir.6.3.e34}.

\end{thebibliography}

\end{document}
